% !TEX root = ../main.tex
\section{Data Description and Preprocessing}
\label{sec:data}

\subsection{Data Sources and Description}

本文使用题面统计及参考资料提供的 Ebola Tracker 病例数据库副本\cite{eboladatabase}。数据覆盖 2026 年 5 月 14 日至 11 月 17 日，共 9\,863 条记录，涉及 8 个省、68 个卫生区。其中，国家级记录 183 条、省级记录 994 条、卫生区级记录 8\,686 条。全国病例序列取自国家级记录，省级与卫生区级记录用于地区比较，三个层级不重复累加。题面统计截至 10 月下旬，与本文数据的截止日期不同。

研究起点已有确诊病例 8 例、死亡 4 人；截至末次报告，累计确诊 11\,613 例、死亡 5\,678 人。因此，研究期内新增确诊为 11\,605 例，新增死亡为 5\,674 人，二者分别与新增字段的合计一致。

\subsection{Data Cleaning}

按日期、层级和地区检查后，未发现重复记录。缺失值不填零，也不对缺报日期插值；一次覆盖多日的新增记录仍归入报告日。累计病例字段完整，但国家级社区死亡和检测样本字段的缺失率分别为 34.97\% 和 26.78\%，卫生区级检测与随访字段则全部缺失，故不据此估计区级检测能力或随访情况。

对累计数的历史下调保留原记录，不强行修成单调序列。比较各区末期规模时，统一使用截止日的累计值，而非历史峰值。所附周表的累计死亡列误用了累计病例数据，本文改由国家级明细重建累计量，起始病例与期间新增分开计算。

死亡周序列还存在一项待核对差异：按明细报告日期汇总，第 25、26 周新增死亡分别为 241、238 人，所附周表则为 346、133 人。两者总量相同，周次分配相差 105 人。本文暂按明细报告日汇总；使用原周表得到的死亡拟合结果仍需复核。

\subsection{Feature Processing}

以 5 月 14 日为起点，每 7 天划分一个左开右闭区间，将国家级新增记录按报告日期求和，得到 27 周序列。第 $n$ 周的新增病例记为 $y_n$，各周合计为 11\,605 例。最后一个区间截至 11 月 19 日，而数据只到 11 月 17 日，因此末周并不完整。

第 10 周报告新增 792 例，其中 7 月 22 日有 369 例被标记为补报，占该周的 46.6\%。本文保留这一周并注明补报情况，不将全部峰值解释为同期传播加快。周序列描述的是报告时间分布，不等同于真实发病时间分布。

病死率由同一层级、同一日期的累计死亡与累计病例计算；随访率仅在分子、分母统计期一致且记录可用时计算。回测按时间划分为前 20 周拟合、后 7 周检验，参数选择不使用验证期数据。病例资料不含逐户接触关系或实际人口流动记录，这些信息需由其他资料或模型假设补充。
